\scanpage{200}
\diagram{\begin{tikzcd}[ampersand replacement=\&, column sep=large,row sep=large]
A \arrow[rr] \& {} \& B \arrow[d,dash,"3"]\\
A^{\mathfrak S_3^+} \arrow[u,"3"] \arrow[rr,dash,"2"] \& {} \& B_1=B^{\mathfrak S_3^+} \arrow[dd,dash,"2"]\\
{} \& B^{\mathfrak S_3}\arrow[ur,dash] \& {}\\
A^{\mathfrak S_3} \arrow[uu,"2"] \arrow[ur,dash] \arrow[rr,dash,"2"] \& {} \& B_0
\end{tikzcd}}
\[
\begin{aligned}
&\alpha+\beta\tau+\alpha'+\beta'\tau'\\
&=(\alpha+\alpha')+(\beta\tau+\beta'(\sigma-\tau))\\
&=\underbrace{\alpha+\alpha'}+\underbrace{\beta'\sigma}+(\beta-\beta')\tau
\end{aligned}
\]
\[\alpha+\alpha'+\beta\sigma-(\alpha+\alpha'+\beta'\sigma)
 =(\beta-\beta')\sigma\qquad\text{OK}\]

\mspage{L 22}
Tout ce que \unsure{je} croyais avoir à dire, à un moment ou à un autre\deleted{\unsure{d}}
de la réflexion, \inserted{(}que ce soit sur moi, ou sur tel autre, ou sur un milieu,
ou sur une époque, je l'ai dit, sans\inserted{jamais} avoir à bousculer mes réticences.
Il a suffi à chaque fois que je les examine avec attention, pour qu'elles s'évanouissent\deleted{\unsure{.}},
sans plus laisser de traces.
\margin{199}
